\documentclass[10pt,a4paper]{article}
\usepackage[margin=16mm,headsep=5mm,footskip=8mm]{geometry}
\usepackage{amsmath,amssymb,booktabs,tabularx,array,graphicx,xcolor,enumitem,fancyhdr,hyperref,float,longtable}
\definecolor{navy}{HTML}{16324F}\definecolor{teal}{HTML}{0F6B67}\definecolor{pale}{HTML}{EEF6F5}\definecolor{warn}{HTML}{FFF4D8}
\hypersetup{colorlinks=true,linkcolor=navy,urlcolor=teal,pdftitle={Thermodynamics Past Exam Booklet Weeks 1-6}}
\setlength{\parindent}{0pt}\setlength{\parskip}{3pt}
\setlist[itemize]{nosep,leftmargin=5mm}\setlist[enumerate]{nosep,leftmargin=6mm}
\pagestyle{fancy}\fancyhf{}\lhead{Thermodynamics Past Exams}\rhead{Weeks 1--6}\cfoot{\thepage}
\newcommand{\questionfigure}[2]{\begin{figure}[H]\centering\IfFileExists{question-images/#1}{\includegraphics[width=0.84\linewidth,height=0.45\textheight,keepaspectratio]{question-images/#1}}{\fbox{\parbox{0.8\linewidth}{Missing source figure: #1}}}\caption{#2}\end{figure}}
\begin{document}
\begin{titlepage}\centering\vspace*{25mm}{\Huge\bfseries Thermodynamics\\Past Exam Booklet\par}\vspace{5mm}{\Large Weeks 1--6 Questions and Clear Answers\par}\vspace{4mm}{\large Mid-semester papers plus selected end-of-year applications\par}\vfill
\fbox{\parbox{0.88\textwidth}{\textbf{How to use this booklet.} Attempt each question before reading its answer. Every solution identifies the system, assumptions, property model, sequential working, units, final result and common trap. Repeated questions are retained because repetition reveals what examiners value. End-of-year questions are clearly marked as stretch applications.}}
\vfill{\small Prepared from the authorised local Thermodynamics assessment files. No live Canvas session was used.\par}\end{titlepage}
\tableofcontents\clearpage

\section{Coverage map and source boundary}
This booklet contains \textbf{22 rewritten past-exam questions}: nineteen from five mid-semester papers and three selected end-of-year refrigeration questions that directly extend Week 5--6 principles.

\begin{center}
\begin{tabular}{p{0.22\linewidth}p{0.13\linewidth}p{0.55\linewidth}}
\toprule
Source & Questions & Main coverage\\\midrule
Test 1 (2016) & 4 & phase diagrams, R-134a piston-cylinder, rigid tank, water tables\\
Test 2 (2018) & 4 & ideal-gas validity, R-134a tables, piston stops, electrical/paddle work\\
Test 3 (2019) & 4 & phase/table skills, electrical heating, ideal-gas process cycle\\
Test 4 (2017 alt.) & 4 & repeated phase skill, steam piston, room heating, R-134a tables\\
2025 test & 3 & sensible heat, spring-loaded piston, condenser heat exchanger\\
End-year selection & 3 & Week 5/6 refrigeration balances and COP as stretch applications\\\bottomrule
\end{tabular}
\end{center}

\begin{center}\fbox{\parbox{0.92\linewidth}{\textbf{Evidence rule.} Official handwritten answers are followed where internally consistent. If a source answer is ambiguous or arithmetically inconsistent, this booklet says so and gives a transparent reconstruction. Property values remain tied to the edition of the course tables used by the paper, so small last-digit differences may occur with another table edition.}}\end{center}

\textbf{What was excluded from the end-of-year set.} Otto, Diesel, detailed Brayton-regenerator and Rankine-cycle problems were not included because they depend on later cycle theory rather than only the first six weeks. The selected refrigeration questions use Week 5 steady-flow balances and Week 6 COP concepts, but are labelled as stretch work because they also require vapour-compression-cycle conventions.

\clearpage\addcontentsline{toc}{section}{Test 1 (2016)}\section*{Test 1}

\subsection*{Q1}
\textbf{Rewritten question.} For the unlabeled $P$--$v$ diagram of a pure substance, identify the phase regions, key features, and process lines requested in Figure 1.

\textbf{Source page/marks.} source/test1.pdf, p.~2; 5 marks.

\textbf{Vague-data note.} This is a diagram-labelling question, so the answer is taken from the handwritten labelled solution page rather than from a numerical calculation.

\textbf{System and boundary.} Not a calculation problem; a conceptual pure-substance phase map on a $P$--$v$ diagram.

\textbf{Assumptions.}
\begin{itemize}
  \item The sketch is the standard pure-substance $P$--$v$ phase diagram.
  \item The handwritten labels on the solution page are the official intended identifications.
  \item The constant-temperature boiling path runs from compressed liquid through the two-phase dome to superheated vapour.
\end{itemize}

\textbf{Given / Find.}
\begin{itemize}
  \item Given: an unlabeled $P$--$v$ diagram for a pure substance.
  \item Find: liquid, solid, vapour, liquid+solid, solid+vapour, liquid+vapour, critical point, triple line, constant-pressure sublimation path, and constant-temperature boiling path.
\end{itemize}

\textbf{Governing equations.} Qualitative phase-diagram topology only; no numerical equations are required.

\textbf{Sequential working.}
\begin{itemize}
  \item Left region: solid.
  \item Between the left boundary and the left side of the triple region: liquid+solid.
  \item Inside the dome: liquid+vapour.
  \item Right of the dome: vapour.
  \item Below the triple-point horizontal segment: solid+vapour.
  \item The apex of the dome: critical point.
  \item The horizontal coexistence line at low pressure: triple line.
  \item Constant-pressure sublimation follows the solid--vapour boundary.
  \item Constant-temperature boiling follows the line from compressed liquid to superheated vapour through the two-phase region.
\end{itemize}

\textbf{Property-table values.} None.

\textbf{Boxed final answer.}
\[
\boxed{\begin{gathered}\text{Six phase regions + critical point + triple line}\\\text{sublimation path + boiling path}\end{gathered}}
\]

\textbf{Exam trap.} Do not place the triple line on the dome boundary; it is the low-pressure horizontal coexistence line.

\questionfigure{test1-q1.png}{Official Q1 $P$--$v$ diagram with labels.}

\subsection*{Q2}
\textbf{Rewritten question.} A piston-cylinder device contains 3 kg of R-134a at 200 kPa and quality $x=0.2$. Heat is added until the pressure reaches 400 kPa, after which the piston moves at constant pressure until the volume becomes twice the initial volume. Determine the initial volume, the boundary work, and the total heat transfer.

\textbf{Source page/marks.} source/test1.pdf, pp.~3--4; 16 marks.

\textbf{Vague-data note.} The source diagram is only partially informative, so the process description is taken from the printed text and verified against the handwritten solution: constant-volume heating to 400 kPa, then constant-pressure expansion to $V_3=2V_1$.

\textbf{System and boundary.} Closed system; the refrigerant is the control mass inside the piston-cylinder device. The boundary is the moving piston face and cylinder walls.

\textbf{Assumptions.}
\begin{itemize}
  \item R-134a properties are read from the saturation tables used in the handwritten solution.
  \item The process from 1 to 2 occurs at constant volume while the piston is on the stops, so $W_{1\to2}=0$.
  \item The process from 2 to 3 occurs at constant pressure $400\,\text{kPa}$.
\end{itemize}

\textbf{Given / Find.}
\begin{itemize}
  \item Given: $m=3\,\text{kg}$, $P_1=200\,\text{kPa}$, $x_1=0.2$, $P_3=400\,\text{kPa}$, $V_3=2V_1$.
  \item Find: $V_1$, $W_{1\to3}$, and $Q_{1\to3}$.
\end{itemize}

\textbf{Governing equations.}
\[
V_1=mv_1,\qquad v_1=v_f+x_1v_{fg}
\]
\[
W_{1\to2}=0,\qquad W_{2\to3}=P_3\,(V_3-V_2)
\]
\[
Q_{1\to3}-W_{1\to3}=\Delta U=m(u_3-u_1)
\]

\textbf{Property-table values.}
\begin{itemize}
  \item At 200 kPa: $v_f=0.0007532\,\text{m}^3/\text{kg}$, $v_g=0.09951\,\text{m}^3/\text{kg}$, $u_f=38.26\,\text{kJ/kg}$, $u_{fg}=186.25\,\text{kJ/kg}$.
  \item At 400 kPa: $v_f=0.0007905\,\text{m}^3/\text{kg}$, $v_g=0.051266\,\text{m}^3/\text{kg}$, $u_f=63.61\,\text{kJ/kg}$, $u_{fg}=171.49\,\text{kJ/kg}$.
\end{itemize}

\textbf{Sequential working.}
\[
v_1=0.0007532+0.2(0.09951-0.0007532)=0.020593\,\text{m}^3/\text{kg}
\]
\[
V_1=3(0.020593)=0.06178\,\text{m}^3
\]
\[
V_2=V_1=0.06178\,\text{m}^3,\qquad V_3=2V_1=0.12356\,\text{m}^3
\]
\[
W_{1\to3}=W_{2\to3}=400\times10^3(0.12356-0.06178)=24.711\,\text{kJ}
\]
\[
u_1=38.26+0.2(186.25)=75.51\,\text{kJ/kg}
\]
\[
v_3=\frac{V_3}{m}=\frac{0.12356}{3}=0.041186\,\text{m}^3/\text{kg}
\]
\[
x_3=\frac{v_3-v_f}{v_{fg}}=\frac{0.041186-0.0007905}{0.051266-0.0007905}=0.800
\]
\[
u_3=63.61+0.8(171.49)=200.8\,\text{kJ/kg}
\]
\[
Q_{1\to3}=W_{1\to3}+m(u_3-u_1)=24.711+3(200.802-75.51)=400.6\,\text{kJ}
\]

\textbf{Boxed final answer.}
\[
\boxed{V_1=0.06178\,\text{m}^3,\qquad W_{1\to3}=24.711\,\text{kJ},\qquad Q_{1\to3}=400.6\,\text{kJ}}
\]

\textbf{Exam trap.} The first leg is constant volume, so its boundary work is zero even though heat is added.

\questionfigure{test1-q2.png}{Official Q2 piston-cylinder apparatus and piecewise $P$--$v$ process working.}

\subsection*{Q3}
\textbf{Rewritten question.} A rigid tank is divided into two equal parts by a partition. Initially one side contains 5 kg of water at 200 kPa and 25$^\circ$C, and the other side is evacuated. The partition is removed and the water expands to fill the entire tank, then exchanges heat with the surroundings until the temperature returns to 25$^\circ$C. Determine the tank volume, the final pressure, and the heat transfer.

\textbf{Source page/marks.} source/test1.pdf, p.~5; 16 marks.

\textbf{Vague-data note.} The handwritten solution uses the 25$^\circ$C saturated-water properties and treats the initial state as compressed liquid; the tank is rigid and the final temperature is forced back to 25$^\circ$C.

\textbf{System and boundary.} Closed system; the water in the rigid tank is the control mass. The boundary is fixed, so there is no boundary motion.

\textbf{Assumptions.}
\begin{itemize}
  \item Water at 25$^\circ$C and 200 kPa is a compressed liquid, so $v_1\approx v_f@25^\circ$C.
  \item The tank is rigid, so the total volume is constant and $W_b=0$.
  \item Final state is at 25$^\circ$C and the same specific volume set by the expanded rigid tank.
\end{itemize}

\textbf{Given / Find.}
\begin{itemize}
  \item Given: $m=5\,\text{kg}$, $P_1=200\,\text{kPa}$, $T_1=25^\circ\text{C}$, initially one half evacuated.
  \item Find: $V_{tank}$, $P_2$, and $Q_{1\to2}$.
\end{itemize}

\textbf{Governing equations.}
\[
v_1\approx v_f@25^\circ\text{C},\qquad V_1=mv_1,\qquad V_{tank}=2V_1
\]
\[
v_2=\frac{V_{tank}}{m},\qquad Q_{1\to2}=\Delta U=m(u_2-u_1)
\]

\textbf{Property-table values.}
\begin{itemize}
  \item At 25$^\circ$C: $v_f=0.001003\,\text{m}^3/\text{kg}$, $v_g=43.34\,\text{m}^3/\text{kg}$, $u_f=104.83\,\text{kJ/kg}$, $u_{fg}=2304.3\,\text{kJ/kg}$.
  \item Saturation pressure at 25$^\circ$C: $P_{sat}=3.1698\,\text{kPa}$.
\end{itemize}

\textbf{Sequential working.}
\[
v_1\approx 0.001003\,\text{m}^3/\text{kg}
\]
\[
V_1=5(0.001003)=0.005015\,\text{m}^3
\]
\[
V_{tank}=2V_1=0.01003\,\text{m}^3
\]
\[
v_2=\frac{0.01003}{5}=0.002006\,\text{m}^3/\text{kg}
\]
\[
x_2=\frac{v_2-v_f}{v_{fg}}=\frac{0.002006-0.001003}{43.34-0.001003}=2.314\times10^{-5}
\]
\[
u_1=104.83\,\text{kJ/kg}
\]
\[
u_2=u_f+x_2u_{fg}=104.83+(2.314\times10^{-5})(2304.3)=104.88\,\text{kJ/kg}
\]
\[
P_2=P_{sat}@25^\circ\text{C}=3.1698\,\text{kPa}
\]
\[
Q_{1\to2}=5(104.88-104.83)=0.267\,\text{kJ}
\]

\textbf{Boxed final answer.}
\[
\boxed{V_{tank}=0.01003\,\text{m}^3,\qquad P_2=3.1698\,\text{kPa},\qquad Q_{1\to2}=0.267\,\text{kJ}}
\]

\textbf{Exam trap.} The tank is rigid, so the pressure change comes from heat transfer and phase change, not from boundary work.

\questionfigure{test1-q3.png}{Rigid-tank diagram used in Q3.}

\subsection*{Q4}
\textbf{Rewritten question.} Complete the missing values in the water-steam property table for cases a--f.

\textbf{Source page/marks.} source/test1.pdf, p.~7; 9 marks.

\textbf{Vague-data note.} The printed table is clear and the handwritten page provides the completed values. Row (e) is a superheated-vapour state that was checked against the handwritten interpolation.

\textbf{System and boundary.} Property-table lookup only; no control volume or control mass boundary is needed.

\textbf{Assumptions.}
\begin{itemize}
  \item Steam tables for water are used.
  \item $x=0$ means saturated liquid, $x=1$ means saturated vapour, and an omitted $x$ means it is not applicable.
  \item When $u$ lies between $u_f$ and $u_g$, the state is a saturated mixture and $u=u_f+xu_{fg}$.
\end{itemize}

\textbf{Given / Find.}
\begin{itemize}
  \item Given: the partial table entries for rows a--f.
  \item Find: every missing $T$, $p$, $u$, $x$, and phase description entry.
\end{itemize}

\textbf{Governing equations.}
\[
x=\frac{u-u_f}{u_{fg}},\qquad u=u_f+xu_{fg}
\]
\[
P=P_{sat}(T),\qquad T=T_{sat}(P)
\]

\textbf{Property-table values.}
\begin{itemize}
  \item Row (a), 145$^\circ$C: $P_{sat}=415.68\,\text{kPa}$, $u_f=610.19\,\text{kJ/kg}$, $u_{fg}=1944.2\,\text{kJ/kg}$.
  \item Row (b), 800 kPa: $T_{sat}=170.41^\circ\text{C}$, but $u=209.33\,\text{kJ/kg}<u_f$, so the state is compressed liquid; the handwritten completion gives $T=50^\circ$C.
  \item Row (c), 270$^\circ$C: $P_{sat}=5503.0\,\text{kPa}$ and $u_g=2593.7\,\text{kJ/kg}$.
  \item Row (d), 30 kPa: $T_{sat}=69.09^\circ\text{C}$, $u_f=289.24\,\text{kJ/kg}$, $u_{fg}=2178.5\,\text{kJ/kg}$.
  \item Row (e), 1600 kPa and $u=3120.1\,\text{kJ/kg}$: superheated vapour; handwritten completion gives $T=500^\circ$C.
  \item Row (f), 85$^\circ$C: $P_{sat}=57.868\,\text{kPa}$, $u_f=355.96\,\text{kJ/kg}$.
\end{itemize}

\textbf{Sequential working.}
\begin{itemize}
  \item (a) $x=(2000-610.19)/1944.2=0.7148$; $P=415.68\,\text{kPa}$; saturated mixture.
  \item (b) At 800 kPa, $u< u_f$, so the state is compressed liquid; the completed table gives $T=50^\circ$C.
  \item (c) Saturated vapour at 270$^\circ$C gives $x=1$ and $u=2593.7\,\text{kJ/kg}$.
  \item (d) $x=0.4$ gives $u=289.24+0.4(2178.5)=1160.64\,\text{kJ/kg}$ and $T=69.09^\circ$C.
  \item (e) Interpolating the superheated tables at 1.6 MPa gives $u=3120.1\,\text{kJ/kg}$ at $T=500^\circ$C.
  \item (f) $x=0$ gives saturated liquid at 85$^\circ$C with $P=57.868\,\text{kPa}$ and $u=355.96\,\text{kJ/kg}$.
\end{itemize}

\textbf{Completed final table.}
\begin{center}\small
\begin{tabular}{cccccl}
\toprule
Row & $T(^\circ\mathrm{C})$ & $p(\mathrm{kPa})$ & $u(\mathrm{kJ/kg})$ & $x$ & Phase \\
\midrule
a & 145 & 415.68 & 2000 & 0.7148 & saturated mixture \\
b & 50 & 800 & 209.33 & -- & compressed liquid \\
c & 270 & 5503.0 & 2593.7 & 1.0 & saturated vapour \\
d & 69.09 & 30 & 1160.64 & 0.4 & saturated mixture \\
e & 500 & 1600 & 3120.1 & -- & superheated vapour \\
f & 85 & 57.868 & 355.96 & 0.0 & saturated liquid \\
\bottomrule
\end{tabular}
\end{center}
\[\boxed{\text{All six states classified before table lookup.}}\]

\textbf{Exam trap.} For row (b), do not force a quality value into a compressed-liquid state; $x$ is only defined on the saturation line.

\clearpage\addcontentsline{toc}{section}{Test 2 (2018)}\section*{Test 2}

\subsection*{Q1}
\textbf{Rewritten question.} Under what conditions (temperature, pressure, density) could the ideal gas equation be used to accurately predict properties for steam?

\textbf{Vague/source note.} The source page gives only the qualitative prompt; the handwritten response is a three-point list and no numerical values are required.

\textbf{System and boundary.} Closed-system property statement for steam; no process calculation is needed.

\textbf{Assumptions.}
\begin{itemize}
  \item Steam is sufficiently far from the saturation dome that ideal-gas behavior is a good approximation.
  \item The pressure is low.
  \item The density is low, which corresponds to a large specific volume.
\end{itemize}

\textbf{Given / Find.}
\begin{itemize}
  \item Given: steam, with the condition asked qualitatively in terms of temperature, pressure, and density.
  \item Find: the conditions under which the ideal-gas equation is accurate.
\end{itemize}

\textbf{Governing equations.}
\[
PV=mRT, \qquad Z\approx 1 \text{ for ideal-gas behavior}
\]

\textbf{Property-table values.} None.

\textbf{Sequential working with units.}
\begin{itemize}
  \item High temperature.
  \item Low pressure.
  \item Low density.
\end{itemize}

\textbf{Boxed official answer.}
\[
\boxed{\text{Use the ideal-gas equation for steam at high temperature, low pressure, and low density.}}
\]

\textbf{Exam trap.} Do not say only ``high temperature''; the handwritten answer explicitly includes low pressure and low density as well.

\textbf{Source pages / marks.} Test sheet p. 2 / 3 marks.

\subsection*{Q2}
\textbf{Rewritten question.} Fill in the missing information in the following table for R134a.

\textbf{Vague/source note.} The source page shows a filled table and two worked parts below it. The handwritten values in the table are the official results to preserve.

\textbf{System and boundary.} Property-state lookup for R134a; each row is a separate equilibrium state.

\textbf{Assumptions.}
\begin{itemize}
  \item Each table entry is read from the appropriate saturated, compressed-liquid, or superheated R134a data.
  \item Saturated-mixture rows use the quality relation $h=h_f+xh_{fg}$.
  \item The handwritten values on the source page are taken as the target answers.
\end{itemize}

\textbf{Given / Find.}
\begin{itemize}
  \item Given: the partially completed R134a table below.
  \item Find: the missing $T$, $p$, $h$, $x$, and phase description entries.
\end{itemize}

\textbf{Governing equations.}
\[
 h=h_f+xh_{fg}, \qquad x=\frac{h-h_f}{h_{fg}} \qquad \text{for saturated mixtures}
\]

\textbf{Property-table values.}
\begin{center}
\begin{tabular}{c c c c l}
\hline
& $T\,(^\circ\mathrm{C})$ & $p\,(\mathrm{kPa})$ & $h\,(\mathrm{kJ/kg})$ & $x$ / phase description \\
\hline
a. & $-22$ & $121.72$ & $22.89$ & $0$; saturated liquid \\
b. & $31.31$ & $800$ & $200$ & $0.6$; saturated mixture \\
c. & $-20$ & $132.82$ & $89.36$ & $0.3$; saturated mixture \\
d. & $8$ & $1000$ & $62.69$ & --; compressed liquid \\
e. & $40$ & $200$ & $287.74$ & --; superheated \\
f. & $85$ & $2928.2$ & $279.51$ & $1.0$; saturated vapour \\
\hline
\end{tabular}
\end{center}

\textbf{Sequential working with units.}
\[
\text{b) } x=\frac{h-h_f}{h_{fg}}=\frac{200-95.48}{174.06}=0.6
\]
\[
\text{c) } h=h_f+xh_{fg}=25.47+0.3\times 212.96=89.36\,\mathrm{kJ/kg}
\]

\textbf{Boxed official answer.}
\[
\boxed{\begin{array}{l}
\text{a. } T=-22^\circ\mathrm{C},\; p=121.72\,\mathrm{kPa},\; h=22.89\,\mathrm{kJ/kg},\; x=0,\; \text{saturated liquid} \\
\text{b. } T=31.31^\circ\mathrm{C},\; p=800\,\mathrm{kPa},\; h=200\,\mathrm{kJ/kg},\; x=0.6,\; \text{saturated mixture} \\
\text{c. } T=-20^\circ\mathrm{C},\; p=132.82\,\mathrm{kPa},\; h=89.36\,\mathrm{kJ/kg},\; x=0.3,\; \text{saturated mixture} \\
\text{d. } T=8^\circ\mathrm{C},\; p=1000\,\mathrm{kPa},\; h=62.69\,\mathrm{kJ/kg},\; x=--,\; \text{compressed liquid} \\
\text{e. } T=40^\circ\mathrm{C},\; p=200\,\mathrm{kPa},\; h=287.74\,\mathrm{kJ/kg},\; x=--,\; \text{superheated} \\
\text{f. } T=85^\circ\mathrm{C},\; p=2928.2\,\mathrm{kPa},\; h=279.51\,\mathrm{kJ/kg},\; x=1.0,\; \text{saturated vapour}
\end{array}}
\]

\textbf{Exam trap.} The source has mixed state types: not every row uses quality; rows d and e are not saturated-mixture states.

\textbf{Source pages / marks.} Test sheet p. 3 / 9 marks.

\subsection*{Q3}
\textbf{Rewritten question.} A piston-cylinder device (Figure 1) with a set of stops on the top contains 2.2 kg of saturated liquid water at 150 kPa. Heat is transferred to the water, causing some of the liquid to evaporate and move the piston upwards (at constant pressure). When the piston reaches the stops, the enclosed volume is 45 l. More heat is transferred until the pressure is doubled. Determine: (a) the amount of saturated liquid, if any, in kg, at the final state; (b) the final temperature; (c) the total work and heat transfer; and (d) show the process on a $p$--$v$ diagram (with respect to saturation lines).

\textbf{Vague/source note.} The handwritten solution uses saturated-water tables at 150 kPa and 300 kPa, then evaluates the final state as a saturated mixture at 300 kPa.

\textbf{System and boundary.} Closed system; moving-boundary piston-cylinder with stops, so boundary work occurs only during the constant-pressure part up to the stops.

\textbf{Assumptions.}
\begin{itemize}
  \item Saturated-water properties are taken from the steam tables.
  \item The piston moves at constant pressure until it hits the stops; after that the volume is fixed.
  \item Kinetic and potential energy changes are negligible.
\end{itemize}

\textbf{Given / Find.}
\begin{itemize}
  \item Given: $m=2.2\,\mathrm{kg}$, initial saturated liquid water at $p_1=150\,\mathrm{kPa}$, stop volume $V_2=45\,\ell=0.045\,\mathrm{m}^3$, and final pressure $p_3=300\,\mathrm{kPa}$.
  \item Find: $m_{f,3}$, $T_3$, $W_{1\to3}$, and $Q_{1\to3}$.
\end{itemize}

\textbf{Governing equations.}
\[
V=mv,\qquad x=\frac{v-v_f}{v_{fg}},\qquad u=u_f+xu_{fg}
\]
\[
W_{1\to3}=W_{1\to2}=p_1(V_2-V_1),\qquad Q_{1\to3}=W_{1\to3}+m(u_3-u_1)
\]

\textbf{Property-table values.}
\[
\begin{gathered}
\text{At }150\,\mathrm{kPa}:\quad u_1=u_f=466.97\,\mathrm{kJ/kg},\quad v_1=v_f=0.001053\,\mathrm{m^3/kg},\\
\text{At }300\,\mathrm{kPa}:\quad u_f=561.11,\ u_{fg}=1982.1\,\mathrm{kJ/kg},\\
v_f=0.001073,\ v_{fg}=0.60582\,\mathrm{m^3/kg},\quad T_{sat}=133.53^\circ\mathrm{C}.
\end{gathered}
\]

\textbf{Sequential working with units.}
\[
V_1=mv_1=(2.2\,\mathrm{kg})(0.001053\,\mathrm{m^3/kg})=0.0023166\,\mathrm{m^3}
\]
\[
V_2=45\,\ell=0.045\,\mathrm{m^3},\qquad p_2=p_1=150\,\mathrm{kPa}
\]
\[
V_3=V_2=0.045\,\mathrm{m^3},\qquad v_3=\frac{V_3}{m}=0.02045\,\mathrm{m^3/kg},\qquad p_3=2p_1=300\,\mathrm{kPa}
\]
\[
v_f<v_3<v_g\text{ at }300\,\mathrm{kPa}\Rightarrow x_3=\frac{v_3-v_f}{v_{fg}}=\frac{0.02045-0.001073}{0.60582-0.001073}=0.032
\]
\[
u_3=u_f+x_3u_{fg}=561.11+0.032\times 1982.1=624.54\,\mathrm{kJ/kg}
\]
\[
m_{f,3}=(1-x_3)m=(1-0.032)(2.2\,\mathrm{kg})=2.12\,\mathrm{kg}
\]
\[
T_3=T_{sat}(300\,\mathrm{kPa})=133.53^\circ\mathrm{C}
\]
\[
W_{1\to3}=W_{1\to2}=p_1(V_2-V_1)=(150\,\mathrm{kPa})(0.045-0.0023166\,\mathrm{m^3})=6.4025\,\mathrm{kJ}
\]
\[
Q_{1\to3}=W_{1\to3}+m(u_3-u_1)=6.4025\,\mathrm{kJ}+(2.2\,\mathrm{kg})(624.54-466.97)\,\mathrm{kJ/kg}=353.1\,\mathrm{kJ}
\]

\textbf{Boxed official answer.}
\[
\boxed{m_{f,3}=2.12\,\mathrm{kg},\qquad T_3=133.53^\circ\mathrm{C},\qquad W_{1\to3}=6.4025\,\mathrm{kJ},\qquad Q_{1\to3}=353.1\,\mathrm{kJ}}
\]

\textbf{Exam trap.} The work stops once the piston hits the stops; the second heating stage changes internal energy only because the volume is fixed.

\textbf{Source pages / marks.} Test sheet pp. 4--5 / 16 marks.

\questionfigure{test2-q3.png}{Piston-cylinder device with stops for Q3.}

\subsection*{Q4}
\textbf{Rewritten question.} An insulated rigid vessel with a volume of 3 m$^3$ is fitted with a resistance heater and a paddle wheel as shown in Figure 2. The vessel is initially filled with air at a temperature of 27$^\circ$C and a pressure of 130 kPa. Now, a current of 4 A from a 230 V supply is passed through the resistance heater while the paddle wheel transfers energy at a rate of 150 J/s to the air. Determine the time taken for the air to reach a temperature of 227$^\circ$C. State any assumptions made and also state any tables used for data.

\textbf{Vague/source note.} The source page and the handwritten solution use the ideal-gas model for air plus temperature-dependent $c_v$ data from Table A-2; a constant-$c_v$ check using Table A-18 is also shown.

\textbf{System and boundary.} Closed rigid vessel; no boundary work, with electrical and shaft work input to the air.

\textbf{Assumptions.}
\begin{itemize}
  \item Air behaves as an ideal gas in the temperature range considered.
  \item The vessel is rigid and insulated, so $Q=0$ and $W_b=0$.
  \item Changes in kinetic and potential energy are negligible.
\end{itemize}

\textbf{Given / Find.}
\begin{itemize}
  \item Given: $P_1=130\,\mathrm{kPa}$, $T_1=27^\circ\mathrm{C}=300\,\mathrm{K}$, $V=3\,\mathrm{m^3}$, $T_2=227^\circ\mathrm{C}=500\,\mathrm{K}$, $I=4\,\mathrm{A}$, $V_{el}=230\,\mathrm{V}$, and paddle-wheel transfer rate $150\,\mathrm{J/s}$ to the air.
  \item Find: time $\Delta t$.
\end{itemize}

\textbf{Governing equations.}
\[
Q-W=\Delta U,\qquad \Delta U=m(u_2-u_1),\qquad m=\frac{PV}{RT}
\]
\[
\dot W_{el}=VI,\qquad \dot W_{sh}=150\,\mathrm{J/s},\qquad \Delta t=\frac{m(u_2-u_1)}{\dot W_{el}+\dot W_{sh}}
\]

\textbf{Property-table values.}
\[
R_{air}=0.287\,\mathrm{kJ/(kg\,K)},\qquad c_{v,air}=0.726\,\mathrm{kJ/(kg\,K)}\;\text{(Table A-2)}
\]
\[
\text{Alternative table values shown on the source page: } u_1=214.07\,\mathrm{kJ/kg}\text{ at }300\,\mathrm{K},\; u_2=359.49\,\mathrm{kJ/kg}\text{ at }500\,\mathrm{K}
\]

\textbf{Sequential working with units.}
\[
m=\frac{PV}{RT_1}=\frac{(130\,\mathrm{kPa})(3\,\mathrm{m^3})}{(0.287\,\mathrm{kJ/(kg\,K)})(300\,\mathrm{K})}=4.53\,\mathrm{kg}
\]
\[
\Delta u=c_{v,air}\Delta T=(0.726\,\mathrm{kJ/(kg\,K)})(500-300\,\mathrm{K})=145.2\,\mathrm{kJ/kg}
\]
\[
\dot W_{el}=-(4\,\mathrm{A})(230\,\mathrm{V})=-920\,\mathrm{J/s},\qquad \dot W_{sh}=-150\,\mathrm{J/s}
\]
\[
\Delta t=\frac{m(u_2-u_1)}{-\dot W_{el}-\dot W_{sh}}=\frac{(4.53\,\mathrm{kg})(145.2\,\mathrm{kJ/kg})}{(920+150)\times 10^{-3}\,\mathrm{kJ/s}}=614.7\,\mathrm{s}
\]
\[
614.7\,\mathrm{s}=10\,\mathrm{min}\;14\,\mathrm{s}
\]

\textbf{Boxed official answer.}
\[
\boxed{\Delta t=614.7\,\mathrm{s}\approx 10\,\mathrm{min}\;14\,\mathrm{s}}
\]

\textbf{Exam trap.} The heater power and paddle-wheel power both enter as energy to the air; do not forget the sign convention when writing the closed-system energy balance.

\textbf{Source pages / marks.} Test sheet pp. 8--9 / 8 marks.

\questionfigure{test2-q4.png}{Insulated rigid vessel with resistance heater and paddle wheel for Q4.}

\clearpage\addcontentsline{toc}{section}{Test 3 (2019)}\section*{Test 3 (2019)}

\subsection*{Q1 --- pure-substance diagram}
\textbf{Rewritten question.} Label the six phase regions, critical point, triple line, constant-pressure sublimation path and constant-temperature boiling path on an unlabeled pure-substance $P$--$v$ diagram.

\textbf{Source page / marks.} Test 3 p.~2; 5 marks.

\textbf{System and boundary.} Conceptual property diagram; no system boundary.

\textbf{Assumptions.}
\begin{itemize}
  \item Standard pure-substance $P$--$v$ topology.
  \item The requested boiling path passes from compressed liquid through the two-phase dome to superheated vapour.
\end{itemize}

\textbf{Given / Find.} Given the blank diagram; identify every phase region and requested path.

\questionfigure{test3-q1.png}{Official labeled phase diagram for Test 3 Q1.}

\textbf{Clear answer.} Solid is far left, liquid is between the fusion line and saturated-liquid boundary, vapour is right of the dome, liquid+vapour is inside the dome, and solid+liquid / solid+vapour occupy the corresponding coexistence regions. The critical point is at the dome apex and the triple line at the low-pressure three-phase line.
\[
\boxed{\text{All six regions, critical point, triple line, sublimation and boiling paths correctly located.}}
\]

\textbf{Exam trap.} The triple line and critical point are different features.

\subsection*{Q2 --- water-steam property table}
\textbf{Rewritten question.} Complete the missing temperature, pressure, internal energy, quality and phase entries for six water-steam states.

\textbf{Source page / marks.} Test 3 p.~3; 9 marks.

\textbf{System and boundary.} Independent equilibrium-state lookups.

\textbf{Assumptions.}
\begin{itemize}
  \item Use the course saturated, superheated and compressed-water tables.
  \item Quality is defined only in the saturated-mixture region.
\end{itemize}

\textbf{Given / Find.} Use each supplied property pair to classify the phase first, then fill all missing values.

\textbf{Core relations.}
\[
x=\frac{u-u_f}{u_{fg}},\qquad u=u_f+xu_{fg}.
\]
Row (a): $x=(2000-610.19)/1944.2=0.715$. Row (d): $u=289.24+0.4(2178.5)=1160.64\,\mathrm{kJ/kg}$.

\textbf{Completed official table.}
\begin{center}
\small
\begin{tabular}{ccccc}
\toprule
Row & $T(^\circ\mathrm{C})$ & $P(\mathrm{kPa})$ & $u(\mathrm{kJ/kg})$ & $x$ / phase\\
\midrule
a & 145 & 415.68 & 2000 & 0.715, saturated mixture\\
b & 50 & 800 & 209.34 & --, compressed liquid\\
c & 270 & 5503.0 & 2593.7 & 1.0, saturated vapour\\
d & 69.09 & 30 & 1160.64 & 0.4, saturated mixture\\
e & 500 & 1600 & 3120.1 & --, superheated vapour\\
f & 85 & 57.868 & 355.36 & 0.0, saturated liquid\\
\bottomrule
\end{tabular}
\end{center}
\[
\boxed{\text{Classify phase before choosing the table and using quality.}}
\]

\textbf{Exam trap.} Do not force quality into rows (b) or (e).

\subsection*{Q3 --- R-134a with heat and electrical work}
\textbf{Rewritten question.} A piston-cylinder contains $2\,\mathrm{kg}$ saturated liquid R-134a and $10\,\mathrm{kg}$ saturated R-134a vapour at $200\,\mathrm{kPa}$. At constant pressure, $250\,\mathrm{kJ}$ of heat is added while a $230\,\mathrm{V}$ source supplies a resistance element for $6\,\mathrm{min}$. If the final temperature is $70^\circ\mathrm{C}$, determine the current and show the process on a $T$--$v$ diagram.

\textbf{Source page / marks.} Test 3 pp.~4--5; 10 marks.

\textbf{System and boundary.} Closed system with moving boundary; heat and electrical work enter the refrigerant.

\textbf{Assumptions.}
\begin{itemize}
  \item Quasi-equilibrium constant-pressure process at $200\,\mathrm{kPa}$.
  \item Negligible KE/PE changes.
  \item R-134a table properties are those used in the official solution.
\end{itemize}

\textbf{Given / Find.} $m_f=2\,\mathrm{kg}$, $m_g=10\,\mathrm{kg}$, $P=200\,\mathrm{kPa}$, $Q=250\,\mathrm{kJ}$, $V_e=230\,\mathrm{V}$, $\Delta t=360\,\mathrm{s}$, $T_2=70^\circ\mathrm{C}$; find $I$ and the process path.

\questionfigure{test3-q3.png}{Piston-cylinder and official $T$--$v$ path for Test 3 Q3.}

\textbf{Sequential working.} Initial quality is the vapour mass fraction:
\[
x_1=\frac{10}{12}=0.833
\]
At $200\,\mathrm{kPa}$, $h_f=38.41$ and $h_{fg}=206.09\,\mathrm{kJ/kg}$:
\[
h_1=h_f+x_1h_{fg}=38.41+0.833(206.09)=210.15\,\mathrm{kJ/kg}
\]
At $200\,\mathrm{kPa}$ and $70^\circ\mathrm{C}$, the final state is superheated and $h_2=315.03\,\mathrm{kJ/kg}$.

For a constant-pressure closed-system process, combining boundary work with $\Delta U$ gives an enthalpy balance:
\[
Q+W_{e,in}=m(h_2-h_1)
\]
\[
W_{e,in}=12(315.03-210.15)-250=1008.56\,\mathrm{kJ}
\]
\[
W_{e,in}=V_e I\Delta t\quad\Rightarrow\quad I=\frac{1008.56\times10^3}{230(360)}=12.18\,\mathrm{A}
\]
\[
\boxed{I=12.18\,\mathrm{A}}
\]

\textbf{Exam trap.} Electrical work is into the system. At constant pressure it is simplest to write the result using $\Delta H$, which already accounts for boundary work.

\subsection*{Q4 --- three-process ideal-gas cycle}
\textbf{Rewritten question.} A piston-cylinder contains $0.2\,\mathrm{kg}$ of air initially at $1.5\,\mathrm{MPa}$ and $200^\circ\mathrm{C}$. It expands at constant pressure until its volume doubles, then cools at constant volume, then is compressed polytropically back to the initial state with $n=1.2$. Determine work for each process and net cycle work, and sketch the $P$--$V$ cycle.

\textbf{Source page / marks.} Test 3 pp.~8--9; 13 marks.

\textbf{System and boundary.} Closed ideal-gas system in a piston-cylinder.

\textbf{Assumptions.}
\begin{itemize}
  \item Air is ideal with $R=0.287\,\mathrm{kJ/(kg\,K)}$.
  \item Quasi-equilibrium boundary work; negligible KE/PE changes.
  \item The return path obeys $PV^{1.2}=\text{constant}$.
\end{itemize}

\textbf{Given / Find.} $m=0.2\,\mathrm{kg}$, $P_1=1500\,\mathrm{kPa}$, $T_1=473\,\mathrm{K}$, $V_2=2V_1$, $V_3=V_2$, $n=1.2$; find $W_{12}$, $W_{23}$, $W_{31}$ and $W_{net}$.

\questionfigure{test3-q4.png}{Official $P$--$V$ cycle sketch for Test 3 Q4.}

\textbf{Sequential working.}
\[
V_1=\frac{mRT_1}{P_1}=\frac{0.2(0.287)(473)}{1500}=0.0181\,\mathrm{m^3},\qquad V_2=V_3=0.0362\,\mathrm{m^3}
\]
\[
W_{12}=P_1(V_2-V_1)=1500(0.0362-0.0181)=27.2\,\mathrm{kJ}
\]
\[
W_{23}=0
\]
Since $P_3V_3^n=P_1V_1^n$,
\[
P_3=P_1\left(\frac{V_1}{V_3}\right)^n=1500(1/2)^{1.2}=652.9\,\mathrm{kPa}
\]
\[
W_{31}=\frac{P_1V_1-P_3V_3}{1-n}=-17.6\,\mathrm{kJ}
\]
\[
W_{net}=27.2+0-17.6=9.6\,\mathrm{kJ}
\]
\[
\boxed{W_{12}=27.2\,\mathrm{kJ},\quad W_{23}=0,\quad W_{31}=-17.6\,\mathrm{kJ},\quad W_{net}=9.6\,\mathrm{kJ}}
\]

\textbf{Exam trap.} Compression work is negative under the work-out convention; the net positive area means the cycle produces work.

\clearpage\addcontentsline{toc}{section}{Test 4 (2017 alternative)}\section*{Test 4 (2017 alternative paper)}

\subsection*{Q1 --- pure-substance $P$--$v$ diagram}
\textbf{Rewritten question.} On the unlabeled $P$--$v$ diagram, identify liquid, solid, vapour, liquid+solid, solid+vapour and liquid+vapour regions; mark the critical point and triple line; draw a constant-pressure sublimation process and a constant-temperature boiling process from compressed liquid to superheated vapour.

\textbf{Source page / marks.} Test 4 p.~1; 5 marks.

\textbf{Source note.} This question repeats the same high-value diagram skill as Test 1 Q1 and Test 3 Q1.

\textbf{System and boundary.} Conceptual pure-substance state map; no control boundary is required.

\textbf{Assumptions.}
\begin{itemize}
  \item The sketch is the standard $P$--$v$ topology for a substance that contracts on freezing.
  \item In the liquid--vapour region, constant-temperature boiling is also constant pressure.
\end{itemize}

\textbf{Given / Find.} Given an unlabeled diagram; find all regions, features and the two requested paths.

\questionfigure{test4-q1.png}{Fully labeled official $P$--$v$ answer for Test 4 Q1.}

\textbf{Clear answer.}
\begin{itemize}
  \item Solid lies to the far left; liquid lies between the fusion boundary and saturated-liquid line; vapour lies right of the dome.
  \item Liquid+solid lies between the solid and liquid boundaries; solid+vapour lies along/below the low-pressure sublimation region; liquid+vapour lies inside the dome.
  \item Critical point: dome apex. Triple line: low-pressure horizontal three-phase line.
  \item Sublimation: horizontal constant-$P$ path from solid to vapour below the triple point.
  \item Boiling: isothermal path from compressed liquid through the dome and into superheated vapour.
\end{itemize}
\[
\boxed{\text{Label all six regions, the critical point, triple line, and both process paths as shown.}}
\]

\textbf{Exam trap.} Do not put the triple line at the critical point.

\subsection*{Q2 --- steam piston-cylinder with stops}
\textbf{Rewritten question.} A piston-cylinder contains $2\,\mathrm{kg}$ of steam at $300\,\mathrm{kPa}$ with quality $x_1=0.25$, while its piston rests on stops. Heat is added at constant volume until $500\,\mathrm{kPa}$, where the piston starts moving. Heating continues at $500\,\mathrm{kPa}$ until the volume is twice its initial value. Determine initial volume, boundary work, total heat transfer, and show the path on a $P$--$v$ diagram.

\textbf{Source page / marks.} Test 4 pp.~2--3; 16 marks.

\textbf{System and boundary.} Closed system; steam bounded by the piston-cylinder walls and moving piston.

\textbf{Assumptions.}
\begin{itemize}
  \item Quasi-equilibrium piston motion; negligible kinetic and potential energy changes.
  \item $1\to2$ is constant volume and $2\to3$ is constant pressure.
  \item Saturated-water tables apply at states 1 and 3.
\end{itemize}

\textbf{Given / Find.} $m=2\,\mathrm{kg}$, $P_1=300\,\mathrm{kPa}$, $x_1=0.25$, $P_2=P_3=500\,\mathrm{kPa}$ and $V_3=2V_1$; find $V_1$, $W$, $Q$, and process path.

\questionfigure{test4-q2.png}{Piston-cylinder apparatus and official $P$--$v$ path for Test 4 Q2.}

\textbf{Property values.}
\[
\begin{array}{ll}
300\,\mathrm{kPa}:&v_f=0.001073,\ v_g=0.60582\,\mathrm{m^3/kg},\ u_f=561.11,\ u_{fg}=1982.1\,\mathrm{kJ/kg},\\
500\,\mathrm{kPa}:&v_f=0.001093,\ v_g=0.37483\,\mathrm{m^3/kg},\ u_f=639.54,\ u_{fg}=1921.2\,\mathrm{kJ/kg}.
\end{array}
\]

\textbf{Sequential working.}
\[
v_1=v_f+x_1(v_g-v_f)=0.001073+0.25(0.60582-0.001073)=0.15226\,\mathrm{m^3/kg}
\]
\[
V_1=mv_1=2(0.15226)=0.30452\,\mathrm{m^3},\quad V_3=2V_1=0.60904\,\mathrm{m^3}
\]
\[
W_{13}=0+P_2(V_3-V_2)=500(0.60904-0.30452)=152.26\,\mathrm{kJ}
\]
\[
u_1=561.11+0.25(1982.1)=1056.64\,\mathrm{kJ/kg}
\]
\[
v_3=V_3/m=0.30452\,\mathrm{m^3/kg},\quad x_3=\frac{0.30452-0.001093}{0.37483-0.001093}=0.812
\]
\[
u_3=639.54+0.812(1921.2)=2199.31\,\mathrm{kJ/kg}
\]
\[
Q_{13}=W_{13}+m(u_3-u_1)=152.26+2(2199.31-1056.64)=2437.6\,\mathrm{kJ}
\]
\[
\boxed{V_1=0.30452\,\mathrm{m^3},\quad W_{out}=152.26\,\mathrm{kJ},\quad Q_{in}=2437.6\,\mathrm{kJ}}
\]

\textbf{Exam trap.} Use total volume in $P\Delta V$, but specific volume in table-quality calculations.

\subsection*{Q3 --- heating a room}
\textbf{Rewritten question.} A $5\times6\times3\,\mathrm{m}$ room at atmospheric pressure is heated by a steam radiator supplying $15{,}000\,\mathrm{kJ/h}$. A $200\,\mathrm{W}$ fan distributes the air, and heat loss is $100\,\mathrm{kJ/min}$. The room is heated from $7^\circ\mathrm{C}$ to $22^\circ\mathrm{C}$. (a) Estimate the heating time by treating room air as a constant-volume system. (b) Discuss whether that assumption is reasonable.

\textbf{Source page / marks.} Test 4 pp.~4--5; 10 marks.

\textbf{System and boundary.} Approximate closed, rigid system containing the room air.

\textbf{Assumptions.}
\begin{itemize}
  \item Ideal-gas air; atmospheric pressure $101.3\,\mathrm{kPa}$ initially.
  \item Uniform room-air temperature; negligible KE/PE changes.
  \item Radiator input, fan work and stated heat loss are constant during heating.
\end{itemize}

\textbf{Given / Find.} $V=90\,\mathrm{m^3}$, $T_1=280\,\mathrm{K}$, $T_2=295\,\mathrm{K}$ and the three energy rates; find $\Delta t$ and judge the model.

\questionfigure{test4-q3.png}{Room, radiator, fan and heat-loss diagram for Test 4 Q3.}

\textbf{Sequential working.}
\[
m=\frac{PV}{RT_1}=\frac{101.3(90)}{0.287(280)}=113.45\,\mathrm{kg}
\]
From ideal-gas Table A-17, $u_1=199.75$ at $280\,\mathrm{K}$ and $u_2=210.49\,\mathrm{kJ/kg}$ at $295\,\mathrm{K}$:
\[
\Delta U=m(u_2-u_1)=113.45(210.49-199.75)=1218.47\,\mathrm{kJ}
\]
Net input rate:
\[
\dot E_{in}=\frac{15000}{3600}+0.200-\frac{100}{60}=2.70\,\mathrm{kJ/s}
\]
\[
\Delta t=\frac{1218.47}{2.70}=451.3\,\mathrm{s}=7.52\,\mathrm{min}
\]
Using average $c_v=0.7175\,\mathrm{kJ/(kg\,K)}$ gives $7.54\,\mathrm{min}$.
\[
\boxed{\Delta t\approx7.52\text{--}7.54\,\mathrm{min}}
\]

\textbf{Assumption discussion.} The geometric room volume is fixed, but a real room is not an airtight closed system; some air escapes as it warms. Constant volume is nevertheless a reasonable exam approximation because the temperature rise and associated leakage are modest. Constant pressure would require a mass-flow model.

\textbf{Exam trap.} Fan work enters the air, while the stated $100\,\mathrm{kJ/min}$ leaves it.

\subsection*{Q4 --- R-134a property table}
\textbf{Rewritten question.} Complete the missing $T$, $P$, $h$, quality and phase entries for six R-134a states.

\textbf{Source page / marks.} Test 4 p.~6; 9 marks.

\textbf{System and boundary.} Independent equilibrium-state lookups; no process boundary.

\textbf{Assumptions.}
\begin{itemize}
  \item Use the course R-134a saturation, superheated-vapour and compressed-liquid tables.
  \item Quality applies only in the two-phase region.
\end{itemize}

\textbf{Given / Find.} Use each supplied property pair to identify phase and fill all missing entries.

\textbf{Working pattern.}
\[
x=\frac{h-h_f}{h_{fg}}\quad\text{for a mixture},\qquad h\approx h_f(T)\quad\text{for a compressed liquid.}
\]
For row (a), $x=(250-102.34)/167.19=0.883$. For row (c), $h=30.67+0.4(210.23)=114.762\,\mathrm{kJ/kg}$.

\textbf{Completed official table.}
\begin{center}
\small
\begin{tabular}{ccccc}
\toprule
Row & $T(^\circ\mathrm{C})$ & $P(\mathrm{kPa})$ & $h(\mathrm{kJ/kg})$ & $x$ / phase\\
\midrule
a & 36 & 912.35 & 250 & 0.88, saturated mixture\\
b & 50 & 500 & 291.98 & --, superheated vapour\\
c & -16 & 157.38 & 114.762 & 0.4, saturated mixture\\
d & 8.91 & 400 & 255.61 & 1.0, saturated vapour\\
e & -32 & 76.71 & 10.09 & 0.0, saturated liquid\\
f & -2 & 900 & 49.15 & --, compressed liquid\\
\bottomrule
\end{tabular}
\end{center}
\[
\boxed{\text{Use phase first, then select the saturation, superheat, or compressed-liquid table.}}
\]

\textbf{Exam trap.} Do not assign a quality to superheated or compressed-liquid states.

\clearpage\addcontentsline{toc}{section}{2025 Mid-Semester Test}\section*{Test 2025}

\subsection*{Q1}
\textbf{Rewritten question.} A 10 kg system was taken from $25\,^\circ\mathrm{C}$ to $75\,^\circ\mathrm{C}$. Given the specific heat capacity of the system is $c_p = 3000\,\mathrm{J/(kg\cdot ^\circ C)}$.\footnote{Source sheet: p.~1, 9 marks total.}

\textbf{Source page / marks.} Question sheet p.~1; Q1 worth 9 marks, split as (a) 4 marks, (b) 3 marks, (c) 2 marks.

\textbf{System and boundary.} Closed system; no mass crosses the boundary. The energy accounting is for the system only.

\textbf{Assumptions.}
\begin{itemize}
  \item The specific heat is constant over the temperature change.
  \item No work interaction is stated, so the heating is interpreted as pure sensible energy addition.
  \item Heat losses are only introduced explicitly in part (c).
\end{itemize}

\textbf{Given / Find.}
\begin{itemize}
  \item Given: $m=10\,\mathrm{kg}$, $T_1=25\,^\circ\mathrm{C}$, $T_2=75\,^\circ\mathrm{C}$, $c_p=3000\,\mathrm{J/(kg\cdot ^\circ C)}$.
  \item Find: (a) $Q$ added, (b) $\dot Q$, (c) heat lost relative to the no-loss case.
\end{itemize}

\textbf{Governing model.}
\[
Q = mc_p\Delta T,\qquad \dot Q = \frac{Q}{\Delta t}
\]

\textbf{Sequential working with units.}
\[
\Delta T = T_2-T_1 = 75-25 = 50\,^\circ\mathrm{C}
\]
\[
Q = (10\,\mathrm{kg})(3000\,\mathrm{J/(kg\cdot ^\circ C)})(50\,^\circ\mathrm{C})
= 1.50\times 10^6\,\mathrm{J} = 1500\,\mathrm{kJ}
\]
\[
\dot Q = \frac{1500\,\mathrm{kJ}}{20\,\mathrm{s}} = 75\,\mathrm{kW}
\]
\[
Q_{\text{ideal}} = (10)(3000)(80-25)=1.65\times 10^6\,\mathrm{J}=1650\,\mathrm{kJ}
\]
\[
Q_{\text{lost}} = Q_{\text{ideal}}-Q_{\text{actual}} = 1650-1500 = 150\,\mathrm{kJ}
\]

\textbf{Boxed answers.}
\[
\boxed{Q=1500\,\mathrm{kJ},\qquad \dot Q=75\,\mathrm{kW},\qquad Q_{\text{lost}}=150\,\mathrm{kJ}}
\]

\textbf{Exam check.} Units reduce correctly: $\mathrm{J/(kg\,^{\circ}C)}\times\mathrm{kg}\times{}^\circ\mathrm{C}=\mathrm{J}$. The loss in part (c) is the difference between the ideal no-loss heating requirement and the actual heating supplied.

\subsection*{Q2}
\textbf{Rewritten question.} 2 kg of water initially at $95\,^\circ\mathrm{C}$ with quality 0.35 occupies a linear spring-loaded piston--cylinder. It is heated until the pressure reaches 800 kPa and the temperature is $275\,^\circ\mathrm{C}$. Determine the initial and final specific volumes and total work.

\textbf{Source page / marks.} Question sheet p.~1; Q2 worth 12 marks, split as (a) 3 marks, (b) 4 marks, (c) 5 marks.

\textbf{System and boundary.} Closed system; piston--cylinder with a moving boundary. The boundary does work only while the piston moves.

\textbf{Assumptions.}
\begin{itemize}
  \item The piston motion is quasi-equilibrium and the spring gives a linear $P$--$V$ path.
  \item The initial state is a saturated mixture at $95\,^\circ\mathrm{C}$, so saturation data apply.
  \item The final state is superheated at 800 kPa and $275\,^\circ\mathrm{C}$, so superheated-steam data apply.
  \item Kinetic and potential energy changes are negligible.
\end{itemize}

\textbf{Given / Find.}
\begin{itemize}
  \item Given: $m=2\,\mathrm{kg}$, $T_1=95\,^\circ\mathrm{C}$, $x_1=0.35$, $P_2=800\,\mathrm{kPa}$, $T_2=275\,^\circ\mathrm{C}$.
  \item Find: $v_1$, $v_2$, and $W_{1\to2}$.
\end{itemize}

\textbf{Governing model.}
\[
 v_1 = v_f + x_1(v_g-v_f),\qquad v_2 = v(P_2,T_2),\qquad W_{1\to2} = \int_1^2 P\,dV
\]
For a linear spring path,
\[
W_{1\to2}=\frac{P_1+P_2}{2}(V_2-V_1)=m\,\frac{P_1+P_2}{2}(v_2-v_1)
\]

\textbf{Source table values.} Using local steam tables for water:
\[
\begin{array}{l}
P_{\text{sat}}(95\,^\circ\mathrm{C}) \approx 84.6\,\mathrm{kPa},\\
v_f(95\,^\circ\mathrm{C}) \approx 0.00104\,\mathrm{m^3/kg},\\
v_g(95\,^\circ\mathrm{C}) \approx 2.01\,\mathrm{m^3/kg},\\
\text{and at } 800\,\mathrm{kPa},\ 275\,^\circ\mathrm{C}:\ v_2\approx 0.307\,\mathrm{m^3/kg}\ \text{(superheated-table lookup; approximate)}.
\end{array}
\]
If a table entry is read as a symbol rather than a fully trusted number, the calculation below is shown with a clearly labelled approximation.

\textbf{Sequential working with units.}
\[
 v_1 = 0.00104 + 0.35(2.01-0.00104) \approx 0.704\,\mathrm{m^3/kg}
\]
\[
 V_1 = mv_1 = (2\,\mathrm{kg})(0.704\,\mathrm{m^3/kg}) = 1.408\,\mathrm{m^3}
\]
\[
 V_2 = mv_2 \approx (2\,\mathrm{kg})(0.307\,\mathrm{m^3/kg}) = 0.614\,\mathrm{m^3}
\]
Because the device is spring-loaded, the pressure varies linearly between states, so using the average pressure gives the same boundary work as the area under the straight line.
\[
W_{1\to2} \approx \frac{84.6+800}{2}\,\mathrm{kPa}\,(0.614-1.408)\,\mathrm{m^3}
\]
\[
W_{1\to2} \approx (442.3\,\mathrm{kPa})(-0.794\,\mathrm{m^3}) \approx -351\,\mathrm{kJ}
\]

\textbf{Boxed answers.}
\[
\boxed{v_1\approx 0.704\,\mathrm{m^3/kg},\qquad v_2\approx 0.307\,\mathrm{m^3/kg},\qquad W_{1\to2}\approx -351\,\mathrm{kJ}}
\]

\textbf{Exam check and source warning.} The computed path is a net compression, so work by the water is negative even though heat is supplied. This conflicts with the prompt's phrase ``work produced'' and suggests that at least one printed state value may be erroneous. Preserve the literal calculation above, but confirm the intended final state with the lecturer before using it as a model answer.

\subsection*{Q3}
\textbf{Rewritten question.} R-134a at 0.8 MPa and $90\,^\circ\mathrm{C}$ is cooled at 0.8 MPa to $30\,^\circ\mathrm{C}$ by air. Air enters at 100 kPa and $22\,^\circ\mathrm{C}$ with a volume flow rate of $400\,\mathrm{m^3/min}$ and leaves at 95 kPa and $57\,^\circ\mathrm{C}$. Determine the inlet/outlet enthalpies of both streams and the refrigerant mass flow rate.

\textbf{Source page / marks.} Question sheet p.~2; Q3 worth 9 marks, split as (a) 2 marks, (b) 2 marks, (c) 1 mark, (d) 1 mark, (e) 3 marks.

\textbf{System and boundary.} Two coupled steady-flow control volumes: one for the R-134a stream and one for the air stream. The condenser wall is the interaction surface between them.

\textbf{Assumptions.}
\begin{itemize}
  \item Steady flow through both streams.
  \item No shaft work and negligible kinetic/potential energy changes.
  \item Air is an ideal gas, so $h=h(T)$ only.
  \item The condenser is adiabatic to the surroundings, so heat lost by the refrigerant equals heat gained by the air.
\end{itemize}

\textbf{Given / Find.}
\begin{itemize}
  \item Given: R-134a at $P_3=0.8\,\mathrm{MPa}$, $T_3=90\,^\circ\mathrm{C}$, $P_4=0.8\,\mathrm{MPa}$, $T_4=30\,^\circ\mathrm{C}$; air at $P_1=100\,\mathrm{kPa}$, $T_1=22\,^\circ\mathrm{C}$, $\dot V_1=400\,\mathrm{m^3/min}$, $P_2=95\,\mathrm{kPa}$, $T_2=57\,^\circ\mathrm{C}$.
  \item Find: $h_1, h_2, h_3, h_4,$ and $\dot m_{\mathrm{R134a}}$.
\end{itemize}

\textbf{Governing model.}
\[
\dot Q_{\mathrm{air}} = \dot m_{\mathrm{air}}(h_4-h_3),\qquad \dot Q_{\mathrm{R}} = \dot m_{\mathrm{R134a}}(h_2-h_1),\qquad \dot Q_{\mathrm{air}} + \dot Q_{\mathrm{R}} = 0
\]
with the air mass flow rate from the ideal-gas relation,
\[
\dot m_{\mathrm{air}} = \rho_1\dot V_1 = \frac{P_1}{RT_1}\dot V_1
\]

\textbf{Source table values.} From the local property tables / worked resource:
\[
\begin{array}{l}
 h_1(\text{R134a},\ 0.8\,\mathrm{MPa},\ 90\,^\circ\mathrm{C}) \approx 416.03\,\mathrm{kJ/kg},\\
 h_2(\text{R134a},\ 0.8\,\mathrm{MPa},\ 30\,^\circ\mathrm{C}) \approx 218.19\,\mathrm{kJ/kg},\\
 h_3(\text{air},\ 22\,^\circ\mathrm{C}) \approx 295.0\,\mathrm{kJ/kg}\ \text{(Table A-17; approximate)},\\
 h_4(\text{air},\ 57\,^\circ\mathrm{C}) \approx 330.8\,\mathrm{kJ/kg}\ \text{(Table A-17; approximate)}.
\end{array}
\]
Hence,
\[
 h_1-h_2 \approx 197.84\,\mathrm{kJ/kg},\qquad h_4-h_3 \approx 35.8\,\mathrm{kJ/kg}
\]

\textbf{Sequential working with units.}
\[
\dot m_{\mathrm{air}} = \frac{(100\,\mathrm{kPa})(400\,\mathrm{m^3/min})}{(0.2870\,\mathrm{kJ/(kg\,K)})(295.15\,\mathrm{K})}
\times \frac{1\,\mathrm{min}}{60\,\mathrm{s}}
\approx 7.92\,\mathrm{kg/s}
\]
\[
\dot Q_{\mathrm{air}} = \dot m_{\mathrm{air}}(h_4-h_3) \approx (7.92\,\mathrm{kg/s})(35.8\,\mathrm{kJ/kg}) \approx 284\,\mathrm{kW}
\]
For the refrigerant side, the heat rejected has the same magnitude:
\[
\dot Q_{\mathrm{R}} = \dot m_{\mathrm{R134a}}(h_2-h_1)
\]
\[
\dot m_{\mathrm{R134a}} = \frac{\dot Q_{\mathrm{air}}}{h_1-h_2} \approx \frac{284\,\mathrm{kW}}{197.84\,\mathrm{kJ/kg}} \approx 1.44\,\mathrm{kg/s}
\]

\textbf{Boxed answers.}
\[
\boxed{h_1\approx 416.03\,\mathrm{kJ/kg},\quad h_2\approx 218.19\,\mathrm{kJ/kg},\quad h_3\approx 295.0\,\mathrm{kJ/kg},\quad h_4\approx 330.8\,\mathrm{kJ/kg},\quad \dot m_{\mathrm{R134a}}\approx 1.44\,\mathrm{kg/s}}
\]

\textbf{Exam check.} The air and refrigerant balances must oppose each other in sign. If the air enthalpy rise is computed from local table values instead of the rounded approximations above, the refrigerant mass flow rate should move slightly but remain of the same order; do not report false precision beyond the table accuracy.

\clearpage\addcontentsline{toc}{section}{Selected end-of-year stretch questions}\section*{Selected end-of-year questions}

\begin{center}
\fbox{\parbox{0.94\linewidth}{\textbf{Scope note.} These are stretch applications of the Week 5 steady-flow balance and Week 6 refrigerator/COP ideas. They also require later vapour-compression-cycle conventions and R-134a property tables. Otto, Diesel, Brayton-regenerator and Rankine questions were excluded because their cycle theory lies beyond Weeks 1--6.}}
\end{center}

\subsection*{End-of-year Problem 10 --- commercial refrigerator}
\textbf{Rewritten question.} A commercial refrigerator using R-134a keeps a refrigerated space at $-25^\circ\mathrm{C}$. Cooling water enters the condenser at $20^\circ\mathrm{C}$ at $0.25\,\mathrm{kg/s}$ and leaves at $25^\circ\mathrm{C}$. Refrigerant enters the condenser at $1.2\,\mathrm{MPa}$ and $70^\circ\mathrm{C}$ and leaves at $42^\circ\mathrm{C}$. It enters the compressor at $60\,\mathrm{kPa}$ and $-20^\circ\mathrm{C}$. The compressor gains $450\,\mathrm{W}$ of heat from the surroundings. Determine (a) quality at the evaporator inlet, (b) refrigeration load and (c) refrigerator COP.

\textbf{Source page / marks.} Previous Test and Exam Problems, p.~6, Problem 10; official solution pp.~18--19.

\textbf{System and boundary.} Steady-flow control volumes around the condenser, compressor, throttle and evaporator.

\textbf{Assumptions.}
\begin{itemize}
  \item Steady operation; negligible kinetic and potential energy changes.
  \item No pressure loss through each heat exchanger.
  \item Throttling is isenthalpic: $h_4=h_3$.
  \item Liquid-water enthalpy is taken from the water table; R-134a properties are taken from the course tables.
\end{itemize}

\textbf{Given / Find.} Given the four refrigerant state conditions, cooling-water data and compressor heat gain; find $x_4$, $\dot Q_L$ and $\mathrm{COP}_R$.

\questionfigure{endyear-p10.png}{Official $T$--$s$ sketch and state numbering for end-of-year Problem 10.}

\textbf{Property values.}
\[
h_1=240.78,\quad h_2=300.63,\quad h_3=h_4=111.28\quad \mathrm{kJ/kg}
\]
At $60\,\mathrm{kPa}$, $h_f=38.87$ and $h_{fg}=223.36\,\mathrm{kJ/kg}$. For the cooling water, $h_{w,in}=83.915$ and $h_{w,out}=104.83\,\mathrm{kJ/kg}$.

\textbf{Sequential working.}
\[
x_4=\frac{h_4-h_f}{h_{fg}}=\frac{111.28-38.87}{223.36}=0.324\approx0.32
\]
Condenser energy balance:
\[
\dot m_w(h_{w,out}-h_{w,in})=\dot m_R(h_2-h_3)
\]
\[
\dot m_R=\frac{0.25(104.83-83.915)}{300.63-111.28}=0.0276\approx0.028\,\mathrm{kg/s}
\]
\[
\dot Q_H=\dot m_R(h_2-h_3)=0.028(300.63-111.28)=5.30\,\mathrm{kW}
\]
For the compressor, heat enters at $0.45\,\mathrm{kW}$:
\[
\dot W_{in}=\dot m_R(h_2-h_1)-\dot Q_{comp,in}
=0.028(300.63-240.78)-0.45=1.23\,\mathrm{kW}
\]
For the complete cycle, $\dot Q_H=\dot Q_L+\dot W_{in}+\dot Q_{comp,in}$:
\[
\dot Q_L=5.30-1.23-0.45=3.62\,\mathrm{kW}
\]
\[
\mathrm{COP}_R=\frac{\dot Q_L}{\dot W_{in}}=\frac{3.62}{1.23}=2.94
\]

\textbf{Boxed official answers.}
\[
\boxed{x_4\approx0.32,\quad \dot m_R\approx0.028\,\mathrm{kg/s},\quad \dot Q_L=3.62\,\mathrm{kW},\quad \dot W_{in}=1.23\,\mathrm{kW},\quad \mathrm{COP}_R=2.94}
\]

\textbf{Exam trap.} The condenser-water heat rate is $\dot Q_H$, not $\dot Q_L$, and the compressor is not adiabatic.

\subsection*{End-of-year Problem 11 --- evaporator pressure and a 250 kW cycle}
\textbf{Rewritten question.} (A) Briefly explain the main factors influencing evaporator-pressure choice in a vapour-compression refrigerator. (B) A $250\,\mathrm{kW}$ R-134a refrigerator receives saturated vapour at $140\,\mathrm{kPa}$ at the compressor inlet and compresses it to $800\,\mathrm{kPa}$ with compressor isentropic efficiency $85\%$. Refrigerant is saturated liquid before throttling. Sketch the cycle, determine quality after throttling, compressor input power, and COP.

\textbf{Source page / marks.} Previous Test and Exam Problems, pp.~6--7, Problem 11; official solution p.~20.

\textbf{System and boundary.} Four steady-flow components forming a simple vapour-compression cycle.

\textbf{Assumptions.}
\begin{itemize}
  \item Steady operation and negligible kinetic/potential energy changes.
  \item Compressor inlet is saturated vapour and throttle inlet is saturated liquid exactly as stated.
  \item Throttle is isenthalpic; compressor is adiabatic with $\eta_c=0.85$.
\end{itemize}

\questionfigure{endyear-p11.png}{Official state sketch and property working for end-of-year Problem 11.}

\textbf{Part A --- clear conceptual answer.}
\begin{itemize}
  \item The refrigerant saturation temperature should normally be about $5$--$10^\circ\mathrm{C}$ below the refrigerated-space temperature so heat has a practical driving temperature difference.
  \item Evaporator pressure should not be below atmospheric pressure where avoidable, because air leakage into the system becomes likely.
  \item A very low evaporator pressure increases compressor pressure ratio, work and discharge temperature, reducing efficiency and reliability.
\end{itemize}

\textbf{Property values from the official solution.}
\[
h_1=236.04,\quad s_1=0.9322,\quad h_3=h_4=95.42\quad \mathrm{kJ/kg}
\]
At $800\,\mathrm{kPa}$ and $s_{2s}=s_1$, interpolation gives $h_{2s}=267.10\,\mathrm{kJ/kg}$.

\textbf{Sequential working.}
After throttling at the $140\,\mathrm{kPa}$ evaporator pressure,
\[
x_4=\frac{h_4-h_f}{h_{fg}}=\frac{95.42-25.77}{210.27}=0.332
\]
Compressor efficiency:
\[
\eta_c=\frac{h_{2s}-h_1}{h_2-h_1}
\quad\Rightarrow\quad
h_2=h_1+\frac{h_{2s}-h_1}{0.85}=272.1\,\mathrm{kJ/kg}
\]
Evaporator load determines refrigerant flow:
\[
250=\dot m_R(h_1-h_4)
\quad\Rightarrow\quad
\dot m_R=\frac{250}{236.04-95.42}=1.78\,\mathrm{kg/s}
\]
\[
\dot W_{in}=\dot m_R(h_2-h_1)=1.78(272.1-236.04)=64.3\,\mathrm{kW}
\]
The official handwritten page reports $74.33\,\mathrm{kW}$ and hence COP $=3.36$; its intermediate mass-flow/enthalpy arithmetic is internally inconsistent with the values printed on that same page. Using the printed properties consistently gives:
\[
\mathrm{COP}_R=\frac{250}{64.3}=3.89
\]

\textbf{Boxed answer with source discrepancy.}
\[
\boxed{x_4\approx0.332,\quad \dot m_R\approx1.78\,\mathrm{kg/s},\quad \dot W_{in}\approx64.3\,\mathrm{kW},\quad \mathrm{COP}_R\approx3.89}
\]
\begin{center}
\fbox{\parbox{0.92\linewidth}{\textbf{Official-page note:} the handwritten sheet boxes $\dot W_{in}=74.33\,\mathrm{kW}$ and $\mathrm{COP}=3.36$. These do not follow from its displayed $h_1$, $h_{2s}$, $h_3$, $250\,\mathrm{kW}$ load and $85\%$ efficiency. The internally consistent reconstruction above is recommended; retain the source values for lecturer comparison.}}
\end{center}

\textbf{Exam trap.} Do not apply compressor efficiency directly to temperature, and do not hide inconsistent source arithmetic.

\subsection*{End-of-year Problem 12 --- condenser heat exchanger}
\textbf{Rewritten question.} An R-134a vapour-compression refrigerator maintains a space at $-30^\circ\mathrm{C}$. Refrigerant enters the condenser at $1.2\,\mathrm{MPa}$ and $65^\circ\mathrm{C}$ and leaves at $42^\circ\mathrm{C}$. It enters the compressor at $60\,\mathrm{kPa}$ and $-34^\circ\mathrm{C}$, while the compressor gains $450\,\mathrm{W}$ from the surroundings. Cooling water enters the condenser at $0.25\,\mathrm{kg/s}$ and $18^\circ\mathrm{C}$ and leaves at $26^\circ\mathrm{C}$. Sketch the cycle; determine refrigeration capacity, compressor input and COP; explain why the compressor inlet is superheated.

\textbf{Source page / marks.} Previous Test and Exam Problems, pp.~7--8, Problem 12; official solution pp.~21--22.

\textbf{System and boundary.} Refrigeration loop plus a separate condenser-water control volume.

\textbf{Assumptions.}
\begin{itemize}
  \item Steady operation; negligible kinetic/potential energy changes and heat loss from the condenser to surroundings.
  \item Condenser and evaporator pressures are $1.2\,\mathrm{MPa}$ and $60\,\mathrm{kPa}$.
  \item Throttling is isenthalpic.
\end{itemize}

\questionfigure{endyear-p12.png}{Official cycle, condenser diagram and $T$--$s$ sketch for end-of-year Problem 12.}

\textbf{Official property values.}
\[
h_1=230.04,\quad h_2=295.18,\quad h_3=h_4=111.25\quad \mathrm{kJ/kg}
\]
For water, $h_{w,in}=75.47$ and $h_{w,out}=109.04\,\mathrm{kJ/kg}$, so $\Delta h_w=33.57\,\mathrm{kJ/kg}$.

\textbf{Sequential working.}
Condenser balance:
\[
\dot m_R(h_2-h_3)=\dot m_w(h_{w,out}-h_{w,in})
\]
\[
\dot m_R=\frac{0.25(109.04-75.47)}{295.18-111.25}=0.0456\,\mathrm{kg/s}
\]
\[
\dot Q_L=\dot m_R(h_1-h_4)=0.0456(230.04-111.25)=5.40\,\mathrm{kW}
\]
With $0.45\,\mathrm{kW}$ entering the compressor,
\[
\dot W_{in}=\dot m_R(h_2-h_1)-0.45
=0.0456(295.18-230.04)-0.45=2.52\,\mathrm{kW}
\]
Using the source's rounded mass flow gives $2.575\,\mathrm{kW}$ and
\[
\mathrm{COP}_R=\frac{5.404}{2.575}=2.10\ \text{(the handwritten page labels this approximately 2.15).}
\]

\textbf{Boxed answers.}
\[
\boxed{\dot m_R\approx0.0456\,\mathrm{kg/s},\quad \dot Q_L\approx5.40\,\mathrm{kW},\quad \dot W_{in}\approx2.52\text{--}2.58\,\mathrm{kW},\quad \mathrm{COP}_R\approx2.10\text{--}2.15}
\]

\textbf{Why superheat the compressor inlet?} A compressor is designed for vapour, not a liquid--vapour mixture. Liquid droplets can damage valves, pistons or blades. Slight superheat ensures dry vapour enters the compressor.

\textbf{Exam trap.} Do not confuse condenser effectiveness with COP, and include the compressor heat gain with the correct sign.


\clearpage
\section{Final exam-use checklist}
\begin{enumerate}
\item Draw the boundary and name the system type before selecting an equation.
\item Build a state table and classify each phase before reading properties.
\item Write the full energy balance before deleting zero terms.
\item Separate total and specific quantities; $PV$ is energy only with compatible units.
\item Mark constant-$V$, constant-$P$, linear-spring and polytropic paths on the diagram.
\item For electrical devices, determine whether $VI\Delta t$ is work into or out of the system.
\item For a heat exchanger, heat lost by one stream equals heat gained by the other if external losses are neglected.
\item For refrigerators, keep $\dot Q_H$, $\dot Q_L$ and $\dot W_{in}$ distinct and state the COP definition.
\item Box a value with units and perform one physical direction/sign check.
\item If the provided source data conflict, state the conflict instead of silently changing them.
\end{enumerate}
\end{document}
